% ID: SEL-D0024-OB-E02
% ESTADO: validada. S025, cotejada con examen y criterios.
\textbf{a)} La continuidad en $1$ exige $a-3=2+b$.
Las pendientes laterales son $a-6$ y $4$. Para que coincidan
hace falta $a-6=4$, de donde $\boxed{a=10,\ b=5}$.

\textbf{b)} Con $a=3$ y $b=-2$ las derivadas en los interiores
son $3-6x$ a la izquierda y $4x$ a la derecha. Por sus signos,
$f$ crece en $(-\infty,1/2]$, decrece en $[1/2,1]$ y crece
en $[1,\infty)$. El último intervalo puede incluir $1$ porque
$f(1)=0$ y $2x^2-2>0$ para $x>1$. En el empalme, ambas expresiones
coinciden, pero sus pendientes laterales son $-3$ y $4$.

Las segundas derivadas son $-6$ en $(-\infty,1)$ y $4$ en $(1,\infty)$.
La gráfica es cóncava hacia abajo antes de $1$ y convexa después.
En $1$ cambia la curvatura en un punto anguloso, sin tangente única.
